% Template adapted from https://github.com/jgm/pandoc-templates/blob/master/default.latex
% To be used with XeLaTex in memoiR
%%%%%%%%%%%%%%%%%%%%%%%%%%%%%%%%%%%%%%%%%%%%%%%%%%%%%%%%%%%%%%%%%%%%%%%%%%%%%%%%%%%%%%%%%

% Options for packages loaded elsewhere
\PassOptionsToPackage{unicode=true}{hyperref}
\PassOptionsToPackage{hyphens}{url}
\PassOptionsToPackage{dvipsnames,svgnames*,x11names*}{xcolor}
% Right to left support


\documentclass[
  10pt,
  italian,
  a4paper,
  extrafontsizes,onecolumn,openright
  ]{memoir}

% Double (or whatever) spacing

% Math
\usepackage{amssymb, amsmath}
% mathspec: arbitrary math fonts
\usepackage{unicode-math}
\defaultfontfeatures{Scale=MatchLowercase}
\defaultfontfeatures[\rmfamily]{Ligatures=TeX,Scale=1}

% Fonts
\usepackage{lmodern}
\usepackage{fontspec}

% Main font
% Specific sanserif font
% Specific monotype font
\setmonofont[Scale=0.85]{Inconsolata}
% Specific math font
% Chinese, Japanese, Corean fonts

% Use upquote for straight quotes in verbatim environments
\usepackage{upquote}
% Use microtype
\usepackage[]{microtype}
\UseMicrotypeSet[protrusion]{basicmath} % disable protrusion for tt fonts

% Verbatim in note

% Color links
\usepackage{xcolor}

% Strikeout

% Necessary for code chunks

% Listings package

% Tables
\usepackage{longtable,booktabs,tabu}
% Fix footnotes in tables (requires footnote package)
\IfFileExists{footnote.sty}{\usepackage{footnote}\makesavenoteenv{longtable}}{}

% Graphics
\usepackage{graphicx,grffile}
\graphicspath{{images/}}
\makeatletter
\def\maxwidth{\ifdim\Gin@nat@width>\linewidth\linewidth\else\Gin@nat@width\fi}
\def\maxheight{\ifdim\Gin@nat@height>\textheight\textheight\else\Gin@nat@height\fi}
\makeatother
% Scale images if necessary, so that they will not overflow the page
% margins by default, and it is still possible to overwrite the defaults
% using explicit options in \includegraphics[width, height, ...]{}
\setkeys{Gin}{width=\maxwidth,height=\maxheight,keepaspectratio}

% Prevent overfull lines
\setlength{\emergencystretch}{3em}  
\providecommand{\tightlist}{%
  \setlength{\itemsep}{0pt}\setlength{\parskip}{0pt}}

% Number sections for memoir (secnumdepth counter is ignored)
\setsecnumdepth{section}

% Set default figure placement to htbp
\makeatletter
\def\fps@figure{htbp}
\makeatother

% Spacing in lists
\usepackage{enumitem}

% Polyglossia
\usepackage{polyglossia}
\setmainlanguage{it}
\setotherlanguage{en-US}

% BibLaTeX
\usepackage[backend=biber,style=authoryear-ibid,isbn=false,backref=true,giveninits=true,uniquename=init,maxcitenames=2,maxbibnames=150,sorting=nyt,sortcites=false,style=apa]{biblatex}
\addbibresource{refs.bib}

% cslreferences environment required by pandoc > 2.7



%%%%%%%%%%%%%%%%%%%%%%%%%%%%%%%%%%%%%%%%%%%%%%%%%%%%%%%%%%
% memoiR format

% Chapter Summary environment 
\usepackage[tikz]{bclogo}
\newenvironment{Summary}
  {\begin{bclogo}[logo=\bctrombone, noborder=true, couleur=lightgray!50]{In breve}\parindent0pt}
  {\end{bclogo}}
% Syntax:
%
%```{block, type='Summary'}
% Deliver message here.
% ```

% scriptsize code 
\let\oldverbatim\verbatim
\def\verbatim{\oldverbatim\scriptsize}
% Applies to code blocks and R code results
% code chunk options size='scriptsize' applies only to R code and results
% if the code chunk sets a different size, \def\verbatim{...} is prioritary for code results 


% Layout
%%%%%%%%%%%%%%%%%%%%%%%%%%%%%%%%%%%%%%%%%%%%%%%%%%%%%%%%%%

% Based on memoir, style companion
\newcommand{\MemoirChapStyle}{daleif1}
\newcommand{\MemoirPageStyle}{Ruled}

% Space between paragraphs
\usepackage{parskip}
  \abnormalparskip{3pt}

% Adjust margin paragraphs vertical position
\usepackage{marginfix}


% Margins
%%%%%%%%%%%%%%%%%%%%%%%%%%%%%%%%%%%%%%%

% allow use of '-',+','/' ans '*' to make simple length computation
\usepackage{calc}

% Full-width figures utilities
\newlength\widthw % full width
\newlength{\rf}
\newcommand*{\definesHSpace}{
  \strictpagecheck % slower but efficient detection of odd/even pages
  \checkoddpage
  \ifoddpage
  \setlength{\rf}{0mm}
  \else
  \setlength{\rf}{\marginparsep+\marginparwidth}
  \fi
}

\makeatletter
% 1" margins for the front matter.
\newcommand*{\SmallMargins}{
  \setlrmarginsandblock{1.5in}{1.5in}{*}
  \setmarginnotes{0.1in}{0.1in}{0.1in}
 \setulmarginsandblock{1.5in}{1in}{*}
  \checkandfixthelayout
  \ch@ngetext
  \clearpage
  \setlength{\widthw}{\textwidth+\marginparsep+\marginparwidth}
  \footnotesatfoot
  \chapterstyle{\MemoirChapStyle}  % Chapter and page styles must be recalled
  \pagestyle{\MemoirPageStyle}
}

% 3" outer margin for the main matter
\newcommand{\LargeMargins}{\SmallMargins}
\makeatother

% Figure captions and footnotes in outer margins


% Main title page with filigrane
%%%%%%%%%%%%%%%%%%%%%%%%%%%%%%%%%%%%%%%%%%%%%%%%%%%%%%%%%%

% Text blocks
\usepackage[absolute,overlay]{textpos}
  \setlength{\TPHorizModule}{1mm}
  \setlength{\TPVertModule}{1mm}

\newcommand{\MainTitlePage}[2]{
  \SmallMargins % Margins
  \pagestyle{empty} % No header/footer
  \textblockorigin{\stockwidth-\paperwidth-\trimedge}{\trimtop} % recto
  \begin{textblock*}{2mm}(\spinemargin/2,\uppermargin/2)
    \rule{1pt}{\paperheight-\uppermargin}
  \end{textblock*}
  \begin{textblock*}{\paperwidth*2/3}(\paperwidth/5, \paperheight/5)
    \flushright
    \begin{Spacing}{3}
      {\fontfamily{qtm}\selectfont\fontsize{45}{45}\selectfont\textsc{\thetitle}}
    \end{Spacing}
  \end{textblock*}
    \begin{textblock*}{\paperwidth*2/3}(\paperwidth/5, \paperheight/2)
    \flushright
    {\fontfamily{qtm}\huge\theauthor}
  \end{textblock*}
    \begin{textblock*}{\paperwidth*2/3}[0, 1](\spinemargin, \uppermargin+\textheight)
    \normalfont\thedate
  \end{textblock*}
  ~\\ % Print a character or the page will not exist
  \newpage
  \textblockorigin{\trimedge}{\trimtop} % verso
  \begin{textblock*}{\textwidth}(\paperwidth-\spinemargin-\textwidth, \uppermargin)
    #1
  \end{textblock*}
  \begin{textblock*}{\textwidth}[0,1](\paperwidth-\spinemargin-\textwidth, \uppermargin+\textheight+\footskip)
    \centering
    \includegraphics[width=\paperwidth/4]{logo}\\ \bigskip
    #2
  \end{textblock*}
  ~\\ % Print a character or the page will not exist
  \newpage
}

% Clear page and open an even one (\clearpage opens an odd one)
\newcommand{\evenpage}{
  \clearpage
  \strictpagecheck % slower but efficient detection of odd/even pages
  \checkoddpage
  \ifoddpage
    \thispagestyle{empty}
    ~\\ % Print a character or the page will not exist
    \newpage
  \else
    % do nothing
  \fi
}


%% PDF title page to insert
%%%%%%%%%%%%%%%%%%%%%%%%%%%%%%%%%%%%%%%%%%%%%%%%%%%%%%%%%%



%% Bibliography
%%%%%%%%%%%%%%%%%%%%%%%%%%%%%%%%%%%%%%%%%%%%%%%%%%%%%%%%%%

\usepackage[strict,autostyle]{csquotes}
% Repeated citation as author-year-title instead of author-title (modification of footcite:note in verbose-inote.cbx)

%% Table of Contents
%%%%%%%%%%%%%%%%%%%%%%%%%%%%%%%%%%%%%%%%%%%%%%%%%%%%%%%%%%

% fix the typesetting of the part number
\renewcommand\partnumberlinebox[2]{#2\ ---\ }


% Fonts
%%%%%%%%%%%%%%%%%%%%%%%%%%%%%%%%%%%%%%%%%%%%%%%%%%%%%%%%%%


% Hyperref comes last
%%%%%%%%%%%%%%%%%%%%%%%%%%%%%%%%%%%%%%%%%%%%%%%%%%%%%%%%%%

\usepackage{hyperref}
\hypersetup{
  pdftitle={Psicometria},
  pdfauthor={Corrado Caudek},
  colorlinks=true,
  linkcolor=Maroon,
  citecolor=Blue,
  urlcolor=Blue,
  breaklinks=true}

% Don't use monospace font for urls
\urlstyle{same}


% Title, author, date from YAML to LaTeX
%%%%%%%%%%%%%%%%%%%%%%%%%%%%%%%%%%%%%%%%%%%%%%%%%%%%%%%%%%

\title{Psicometria}

\author{Corrado Caudek}

\date{2021-11-11}


% Include headers (preamble.tex) here
%%%%%%%%%%%%%%%%%%%%%%%%%%%%%%%%%%%%%%%%%%%%%%%%%%%%%%%%%%
% Add LaTeX code into the preamble of the document here
\hyphenation{bio-di-ver-si-ty sap-lings}


%%%%%%%%%%%%%%%%%%%%%%%%%%%%%%%%%%%%%%%%%%%%%%%%%%%%%%%%%%%%%%%%%%%%%%%%%
% memoiR dalef3 chapter style 
% https://ctan.crest.fr/tex-archive/info/latex-samples/MemoirChapStyles/MemoirChapStyles.pdf
\usepackage{soul}
\definecolor{nicered}{rgb}{.647,.129,.149}
\makeatletter
\newlength\dlf@normtxtw
\setlength\dlf@normtxtw{\textwidth}
\def\myhelvetfont{\def\sfdefault{mdput}}
\newsavebox{\feline@chapter}
\newcommand\feline@chapter@marker[1][4cm]{%
  \sbox\feline@chapter{%
    \resizebox{!}{#1}{\fboxsep=1pt%
	  \colorbox{nicered}{\color{white}\bfseries\sffamily\thechapter}%
	}}%
  \rotatebox{90}{%
    \resizebox{%
	  \heightof{\usebox{\feline@chapter}}+\depthof{\usebox{\feline@chapter}}}%
	{!}{\scshape\so\@chapapp}}\quad%
  \raisebox{\depthof{\usebox{\feline@chapter}}}{\usebox{\feline@chapter}}%
 }
\newcommand\feline@chm[1][4cm]{%
  \sbox\feline@chapter{\feline@chapter@marker[#1]}%
  \makebox[0pt][l]{% aka \rlap
    \makebox[1cm][r]{\usebox\feline@chapter}%
  }}
\makechapterstyle{daleif1}{
  \renewcommand\chapnamefont{\normalfont\Large\scshape\raggedleft\so}
  \renewcommand\chaptitlefont{\normalfont\huge\bfseries\scshape\color{nicered}}
  \renewcommand\chapternamenum{}
  \renewcommand\printchaptername{}
  \renewcommand\printchapternum{\null\hfill\feline@chm[2.5cm]\par}
  \renewcommand\afterchapternum{\par\vskip\midchapskip}
  \renewcommand\printchaptertitle[1]{\chaptitlefont\raggedleft ##1\par}
}
\makeatother

\DeclareMathOperator{\Var}{Var} % Define variance operator
\DeclareMathOperator{\SD}{SD} % Define sd operator
\DeclareMathOperator{\Cov}{Cov} % Define covariance operator
\DeclareMathOperator{\Corr}{Corr} % Define correlation operator
\DeclareMathOperator{\Me}{Me} % Define mediane operator
\DeclareMathOperator{\Mo}{Mo} % Define mode operator
\DeclareMathOperator{\Bin}{Bin} % Define binomial operator
\DeclareMathOperator{\Bernoulli}{Bernoulli} % Define Bernoulli operator
\DeclareMathOperator{\Poi}{Poi} % Define Poisson operator
\DeclareMathOperator{\Uniform}{Uniform} % Define Uniform operator
\DeclareMathOperator{\Cauchy}{Cauchy} % Define Cauchy operator
\DeclareMathOperator{\elpd}{elpd} % Define elpd operator
\DeclareMathOperator{\lppd}{lppd} % Define lppd operator
\DeclareMathOperator{\LOO}{LOO} % Define LOO operator
\DeclareMathOperator{\B}{\mathscr{B}} % Define Bernoulli operator
\newcommand{\R}{\textsf{R}} % Define R programming language symbol
\newcommand{\E}{\mathbb{E}} % Define expected value operator
\newcommand{\Real}{\mathbb{R}} % Define real number operator
\newcommand{\Prob}{\mathscr{P}}
\DeclareMathOperator*{\argmin}{arg\,min} % thin space, limits on side in displays
\DeclareMathOperator*{\argmax}{arg\,max} % thin space, limits on side in displays

\raggedbottom % allow variable (ragged) site heights
\frenchspacing

\usepackage[
  labelfont=bf, 
  font={small, it} 
]{caption} 
\usepackage{upquote} % print correct quotes in verbatim-environments
\usepackage{empheq} 
\usepackage{xfrac}




\usepackage{booktabs}
\usepackage{longtable}
\usepackage{array}
\usepackage{multirow}
\usepackage{wrapfig}
\usepackage{float}
\usepackage{colortbl}
\usepackage{pdflscape}
\usepackage{tabu}
\usepackage{threeparttable}
\usepackage{threeparttablex}
\usepackage[normalem]{ulem}
\usepackage{makecell}
\usepackage{xcolor}


% End of preamble
%%%%%%%%%%%%%%%%%%%%%%%%%%%%%%%%%%%%%%%%%%%%%%%%%%%%%%%%%%


\usepackage{amsthm}
\newtheorem{theorem}{Teorema}[chapter]
\newtheorem{lemma}{Lemma}[chapter]
\newtheorem{corollary}{Corollario}[chapter]
\newtheorem{proposition}{Proposizione}[chapter]
\newtheorem{conjecture}{Congettura}[chapter]
\theoremstyle{definition}
\newtheorem{definition}{Definizione}[chapter]
\theoremstyle{definition}
\newtheorem{example}{Esempio}[chapter]
\theoremstyle{definition}
\newtheorem{exercise}{Exercizio}[chapter]
\theoremstyle{definition}
\newtheorem{hypothesis}{Hypothesis}[chapter]
\theoremstyle{remark}
\newtheorem*{remark}{Osservazione}
\newtheorem*{solution}{Soluzione}
\begin{document}
\frontmatter

% Title page
%%%%%%%%%%%%%%%%%%%%%%%%%%%%%%%%%%%%%%%%%%%%%%%%%%%%%%%%%%


\MainTitlePage{Questo documento è stato realizzato con:

\begin{itemize}
  \item \LaTeX\; e la classe memoir (\url{http://www.ctan.org/pkg/memoir});
  \item $\R$ (\url{http://www.r-project.org/}) e RStudio (\url{http://www.rstudio.com/});
  \item bookdown (\url{http://bookdown.org/}) e memoiR (\url{https://ericmarcon.github.io/memoiR/}).
\end{itemize}}{Nel blog della mia pagina personale sono forniti alcuni approfondimenti degli argomenti qui trattati.

\url{https://ccaudek.github.io/caudeklab/}}


% Before Body
%%%%%%%%%%%%%%%%%%%%%%%%%%%%%%%%%%%%%%%%%%%%%%%%%%%%%%%%%%





% Contents
%%%%%%%%%%%%%%%%%%%%%%%%%%%%%%%%%%%%%%%%%%%%%%%%%%%%%%%%%%

\LargeMargins
{
\hypersetup{linkcolor=}
\setcounter{tocdepth}{2}
\tableofcontents
}


% Body
%%%%%%%%%%%%%%%%%%%%%%%%%%%%%%%%%%%%%%%%%%%%%%%%%%%%%%%%%%

\LargeMargins
\scriptsize

\normalsize

\chapter*{}

\vfill

\scriptsize

\normalsize

\scriptsize

Copyright \(\copyright\) 2022.

\normalsize

Data della versione presente: Novembre 11, 2021.

\hypertarget{prefazione}{%
\chapter{Prefazione}\label{prefazione}}

\textbf{Data Science per psicologi} contiene il materiale delle lezioni dell'insegnamento di \emph{Psicometria B000286} (A.A. 2021/2022) rivolto agli studenti del primo anno del Corso di Laurea in Scienze e Tecniche Psicologiche dell'Università degli Studi di Firenze.

L'insegnamento di Psicometria si propone di fornire agli studenti un'introduzione all'analisi dei dati in psicologia.
Le conoscenze/competenze che verranno sviluppate in questo insegnamento sono quelle della \emph{Data science}, ovvero le conoscenze/competenze che si pongono all'intersezione tra statistica (ovvero, richiedono la capacità di comprendere teoremi statistici) e informatica (ovvero, richiedono la capacità di sapere utilizzare un software).

\hypertarget{la-psicologia-e-la-data-science}{%
\section*{La psicologia e la Data Science}\label{la-psicologia-e-la-data-science}}
\addcontentsline{toc}{section}{La psicologia e la Data Science}

\begin{quote}
It's worth noting, before getting started, that this material is hard. If you find yourself confused at any point, you are normal. Any sense of confusion you feel is just your brain correctly calibrating to the subject matter. Over time, confusion is replaced by comprehension {[}\ldots{]} --- Richard McElreath
\end{quote}

Sembra sensato spendere due parole su un tema che è importante per gli studenti: quello indicato dal titolo di questo Capitolo. È ovvio che agli studenti di psicologia la statistica non piace. Se piacesse, forse studierebbero Data Science e non psicologia; ma non lo fanno. Di conseguenza, gli studenti di psicologia si chiedono: ``perché dobbiamo perdere tanto tempo a studiare queste cose quando in realtà quello che ci interessa è tutt'altro?'\,' Questa è una bella domanda.

C'è una ragione molto semplice che dovrebbe farci capire perché la Data Science è così importante per la psicologia. Infatti, a ben pensarci, la psicologia è una disciplina intrinsecamente statistica, se per statistica intendiamo quella disciplina che studia la variazione delle caratteristiche degli individui nella popolazione. La psicologia studia \emph{gli individui} ed è proprio la variabilità inter- e intra-individuale ciò che vogliamo descrivere e, in certi casi, predire. In questo senso, la psicologia è molto diversa dall'ingegneria, per esempio. Le proprietà di un determinato ponte sotto certe condizioni, ad esempio, sono molto simili a quelle di un altro ponte, sotto le medesime condizioni. Quindi, per un ingegnere la statistica è poco importante: le proprietà dei materiali sono unicamente dipendenti dalla loro composizione e restano costanti. Ma lo stesso non può dirsi degli individui: ogni individuo è unico e cambia nel tempo. E le variazioni tra gli individui, e di un individuo nel tempo, sono l'oggetto di studio proprio della psicologia: è dunque chiaro che i problemi che la psicologia si pone sono molto diversi da quelli affrontati, per esempio, dagli ingegneri. Questa è la ragione per cui abbiamo tanto bisogno della \emph{data science} in psicologia: perché la \emph{data science} ci consente di descrivere la variazione e il cambiamento. E queste sono appunto le caratteristiche di base dei fenomeni psicologici.

Sono sicuro che, leggendo queste righe, a molti studenti sarà venuta in mente la seguente domanda: perché non chiediamo a qualche esperto di fare il ``lavoro sporco'' (ovvero le analisi statistiche) per noi, mentre noi (gli psicologi) ci occupiamo solo di ciò che ci interessa, ovvero dei problemi psicologici slegati dai dettagli ``tecnici'' della \emph{data science}?
La risposta a questa domanda è che non è possibile progettare uno studio psicologico sensato senza avere almeno una comprensione rudimentale della \emph{data science}. Le tematiche della \emph{data science} non possono essere ignorate né dai ricercatori in psicologia né da coloro che svolgono la professione di psicologo al di fuori dell'Università. Infatti, anche i professionisti al di fuori dall'università non possono fare a meno di leggere la letteratura psicologica più recente: il continuo aggiornamento delle conoscenze è infatti richiesto dalla deontologia della professione. Ma per potere fare questo è necessario conoscere un bel po' di \emph{data science}! Basta aprire a caso una rivista specialistica di psicologia per rendersi conto di quanto ciò sia vero: gli articoli che riportano i risultati delle ricerche psicologiche sono zeppi di analisi statistiche e di modelli formali. E la comprensione della letteratura psicologica rappresenta un requisito minimo nel bagaglio professionale dello psicologo.

Le considerazioni precedenti cercano di chiarire il seguente punto: la \emph{data science} non è qualcosa da studiare a malincuore, in un singolo insegnamento universitario, per poi poterla tranquillamente dimenticare. Nel bene e nel male, gli psicologi usano gli strumenti della \emph{data science} in tantissimi ambiti della loro attività professionale: in particolare quando costruiscono, somministrano e interpretano i test psicometrici. È dunque chiaro che possedere delle solide basi di \emph{data science} è un tassello imprescindibile del bagaglio professionale dello psicologo. In questo insegnamento verrano trattati i temi base della \emph{data science} e verrà adottato un punto di vista bayesiano, che corrisponde all'approccio più recente e sempre più diffuso in psicologia.

\hypertarget{come-studiare}{%
\section*{Come studiare}\label{come-studiare}}
\addcontentsline{toc}{section}{Come studiare}

\begin{quote}
I know quite certainly that I myself have no special talent. Curiosity, obsession and dogged endurance, combined with self-criticism, have brought me to my ideas. --- Albert Einstein
\end{quote}

Il giusto metodo di studio per prepararsi all'esame di Psicometria è quello di seguire attivamente le lezioni, assimilare i concetti via via che essi vengono presentati e verificare in autonomia le procedure presentate a lezione. Incoraggio gli studenti a farmi domande per chiarire ciò che non è stato capito appieno. Incoraggio gli studenti a utilizzare i forum attivi su Moodle e, soprattutto, a svolgere gli esercizi proposti su Moodle. I problemi forniti su Moodle rappresentano il livello di difficoltà richiesto per superare l'esame e consentono allo studente di comprendere se le competenze sviluppate fino a quel punto sono sufficienti rispetto alle richieste dell'esame.

La prima fase dello studio, che è sicuramente individuale, è quella in cui è necessario acquisire le conoscenze teoriche relative ai problemi che saranno presentati all'esame. La seconda fase di studio, che può essere facilitata da scambi con altri e da incontri di gruppo, porta ad acquisire la capacità di applicare le conoscenze: è necessario capire come usare un software (\R) per applicare i concetti statistici alla specifica situazione del problema che si vuole risolvere. Le due fasi non sono però separate: il saper fare molto spesso ci aiuta a capire meglio.

\hypertarget{sviluppare-un-metodo-di-studio-efficace}{%
\section*{Sviluppare un metodo di studio efficace}\label{sviluppare-un-metodo-di-studio-efficace}}
\addcontentsline{toc}{section}{Sviluppare un metodo di studio efficace}

\begin{quote}
Memorization is not learning. --- Richard Phillips Feynman
\end{quote}

Avendo insegnato molte volte in passato un corso introduttivo di analisi dei dati ho notato nel corso degli anni che gli studenti con l'atteggiamento mentale che descriverò qui sotto generalmente ottengono ottimi risultati. Alcuni studenti sviluppano naturalmente questo approccio allo studio, ma altri hanno bisogno di fare uno sforzo per maturarlo. Fornisco qui sotto una breve descrizione del ``metodo di studio'\,' che, nella mia esperienza, è il più efficace per affrontare le richieste di questo insegnamento \autocite{burger20125}.

\begin{itemize}
\tightlist
\item
  Dedicate un tempo sufficiente al materiale di base, apparentemente facile; assicuratevi di averlo capito bene. Cercate le lacune nella vostra comprensione. Leggere presentazioni diverse dello stesso materiale (in libri o articoli diversi) può fornire nuove intuizioni.
\end{itemize}

\begin{itemize}
\item
  Gli errori che facciamo sono i nostri migliori maestri. Istintivamente cerchiamo di dimenticare subito i nostri errori. Ma il miglior modo di imparare è apprendere dagli errori che commettiamo. In questo senso, una soluzione corretta è meno utile di una soluzione sbagliata. Quando commettiamo un errore questo ci fornisce un'informazione importante: ci fa capire qual è il materiale di studio sul quale dobbiamo ritornare e che dobbiamo capire meglio.
\item
  C'è ovviamente un aspetto ``psicologico'' nello studio. Quando un esercizio o problema ci sembra incomprensibile, la cosa migliore da fare è dire: ``mi arrendo'', ``non ho idea di cosa fare!''. Questo ci rilassa: ci siamo già arresi, quindi non abbiamo niente da perdere, non dobbiamo più preoccuparci. Ma non dobbiamo fermarci qui. Le cose ``migliori'' che faccio (se ci sono) le faccio quando non ho voglia di lavorare. Alle volte, quando c'è qualcosa che non so fare e non ho idea di come affontare, mi dico: ``oggi non ho proprio voglia di fare fatica'', non ho voglia di mettermi nello stato mentale per cui ``in 10 minuti devo risolvere il problema perché dopo devo fare altre cose''. Però ho voglia di \emph{divertirmi} con quel problema e allora mi dedico a qualche aspetto ``marginale'' del problema, che so come affrontare, oppure considero l'aspetto più difficile del problema, quello che non so come risolvere, ma invece di cercare di risolverlo, guardo come altre persone hanno affrontato problemi simili, opppure lo stesso problema in un altro contesto. Non mi pongo l'obiettivo ``risolvi il problema in 10 minuti'', ma invece quello di farmi un'idea ``generale'' del problema, o quello di capire un caso più specifico e più semplice del problema. Senza nessuna pressione. Infatti, in quel momento ho deciso di non lavorare (ovvero, di non fare fatica). Va benissimo se ``parto per la tangente'', ovvero se mi metto a leggere del materiale che sembra avere poco a che fare con il problema centrale (le nostre intuizioni e la nostra curiosità solitamente ci indirizzano sulla strada giusta). Quando faccio così, molto spesso trovo la soluzione del problema che mi ero posto e, paradossalmente, la trovo in un tempo minore di quello che, in precedenza, avevo dedicato a ``lavorare'' al problema. Allora perché non faccio sempre così? C'è ovviamente l'aspetto dei ``10 minuti'' che non è sempre facile da dimenticare. Sotto pressione, possiamo solo agire in maniera automatica, ovvero possiamo solo applicare qualcosa che già sappiamo fare. Ma se dobbiamo imparare qualcosa di nuovo, la pressione è un impedimento.
\item
  È utile farsi da soli delle domande sugli argomenti trattati, senza limitarsi a cercare di risolvere gli esercizi che vengono assegnati. Quando studio qualcosa mi viene in mente: ``se questo è vero, allora deve succedere quest'altra cosa''. Allora verifico se questo è vero, di solito con una simulazione. Se i risultati della simulazione sono quelli che mi aspetto, allora vuol dire che ho capito. Se i risultati sono diversi da quelli che mi aspettavo, allora mi rendo conto di non avere capito e ritorno indietro a studiare con più attenzione la teoria che pensavo di avere capito -- e ovviamente mi rendo conto che c'era un aspetto che avevo frainteso. Questo tipo di verifica è qualcosa che dobbiamo fare da soli, in prima persona: nessun altro può fare questo al posto nostro.
\item
  Non aspettatevi di capire tutto la prima volta che incontrate un argomento nuovo.\footnote{Ricordatevi inoltre che gli individui tendono a sottostimare la propria capacità di apprendere \autocite{horn2021underestimating}.} È utile farsi una nota mentalmente delle lacune nella vostra comprensione e tornare su di esse in seguito per carcare di colmarle. L'atteggiamento naturale, quando non capiamo i dettagli di qualcosa, è quello di pensare: ``non importa, ho capito in maniera approssimativa questo punto, non devo preoccuparmi del resto''. Ma in realtà non è vero: se la nostra comprensione è superficiale, quando il problema verrà presentato in una nuova forma, non riusciremo a risolverlo. Per cui i dubbi che ci vengono quando studiamo qualcosa sono il nostro alleato più prezioso: ci dicono esattamente quali sono gli aspetti che dobbiamo approfondire per potere migliorare la nostra preparazione.
\item
  È utile sviluppare una visione d'insieme degli argomenti trattati, capire l'obiettivo generale che si vuole raggiungere e avere chiaro il contributo che i vari pezzi di informazione forniscono al raggiungimento di tale obiettivo. Questa organizzazione mentale del materiale di studio facilita la comprensione. È estremamente utile creare degli schemi di ciò che si sta studiando. Non aspettate che sia io a fornirvi un riepilogo di ciò che dovete imparare: sviluppate da soli tali schemi e tali riassunti.
\item
  Tutti noi dobbiamo imparare l'arte di trovare le informazioni, non solo nel caso di questo insegnamento. Quando vi trovate di fronte a qualcosa che non capite, o ottenete un oscuro messaggio di errore da un software, ricordatevi: ``Google is your friend''.
\end{itemize}

\bigskip

Corrado Caudek

\bigskip

Febbraio 2022

\mainmatter

\hypertarget{chapter-prob-congiunta}{%
\chapter{Probabilità congiunta}\label{chapter-prob-congiunta}}

In psicologia e nella vita quotidiana siamo spesso interessati a studiare problemi di probabilità legati al valore congiunto di due o più variabili casuali. Ad esempio, potremmo misurare il QI dei bambini e il loro peso alla nascita, o l'altezza e il peso delle giraffe, o il livello di inquinamento atmosferico e il tasso di malattie respiratorie nelle città, o il numero di amici di Facebook e l'età. Ci potremmo chiedere che relazione esiste tra le coppie di variabili elencate in precedenza.

Per descrivere la relazione tra due variabili casuali è necessario calcolare la \emph{covarianza} e la \emph{correlazione}. Il calcolo di questi due indici richiede la conoscenza della funzione di probabilità congiunta. Obiettivo di questo Capitolo è chiarire cosa si intende per funzione di probabilità congiunta di due variabili casuali \(X\) e \(Y\). Esamineremo qui in dettaglio il caso discreto.

\hypertarget{funzione-di-probabilituxe0-congiunta}{%
\section{Funzione di probabilità congiunta}\label{funzione-di-probabilituxe0-congiunta}}

Dopo aver trattato della distribuzione di probabilità di una variabile casuale, la quale associa ad ogni evento elementare dello spazio campionario uno ed un solo numero reale, è naturale estendere questo concetto al caso di due o più variabili casuali. Iniziamo a descrivere il caso discreto con un esempio.

Consideriamo l'esperimento casuale corrispondente al lancio di tre monete equilibrate. Lo spazio campionario è

\[
\Omega = \{TTT, TTC, TCT, CTT, CCT, CTC, TCC, CCC\}.
\]

\noindent
Dato che i tre lanci sono tra loro indipendenti, non c'è ragione di aspettarsi che uno degli otto risultati possibili dell'esperimento sia più probabile degli altri, dunque possiamo associare a ciascuno degli otto eventi elementari dello spazio campionario la stessa probabilità, ovvero 1/8.

Su tale spazio campionario consideriamo le variabili aleatorie
\(X \in \{0, 1, 2, 3\}\) che conta il numero delle teste nei tre lanci e \(Y \in \{0, 1\}\) che conta il numero delle teste al primo lancio. Indicando con T = `testa' e C = `croce', si ottiene dunque la situazione riportata nella tabella~\ref{tab:tre-monete-distr-cong-1}.

\begin{longtable}[]{@{}cccc@{}}
\caption{\label{tab:tre-monete-distr-cong-1} Spazio campionario dell'esperimento consistente nel lancio di tre monete equilibrate su cui sono state definite le variabili aleatorie \(X\) e \(Y\).}\tabularnewline
\toprule
\(\omega\) & \(X\) & \(Y\) & \(P(\omega)\) \\
\midrule
\endfirsthead
\toprule
\(\omega\) & \(X\) & \(Y\) & \(P(\omega)\) \\
\midrule
\endhead
\(\omega_1\) = TTT & 3 & 1 & 1/8 \\
\(\omega_2\) = TTC & 2 & 1 & 1/8 \\
\(\omega_3\) = TCT & 2 & 1 & 1/8 \\
\(\omega_4\) = CTT & 2 & 0 & 1/8 \\
\(\omega_5\) = CCT & 1 & 0 & 1/8 \\
\(\omega_6\) = CTC & 1 & 0 & 1/8 \\
\(\omega_7\) = TCC & 1 & 1 & 1/8 \\
\(\omega_8\) = CCC & 0 & 0 & 1/8 \\
\bottomrule
\end{longtable}

Ci poniamo il problema di associare un livello di probabilità ad ogni coppia \((x, y)\) definita su \(\Omega\). La coppia \((X = 0, Y = 0)\) si realizza in corrispondenza di un solo evento elementare, ovvero CCC; avrà dunque una probabilità pari a \(P(X=0, Y=0) = P(CCC) = 1/8\). Nel caso della coppia \((X = 1, Y = 0)\) ci sono due eventi elementari che danno luogo al risultato considerato, ovvero, CCT e CTC; la probabilità \(P(X=1, Y=0)\) sarà dunque data dall'unione delle probabilità dei due eventi elementari corrispondenti, cioé \(P(X=1, Y=0) = P(CCT \:\cup\: CTC) = 1/8 + 1/8 = 1/4\). Riportiamo qui sotto i calcoli svolti per tutti i possibili valori di \(X\) e \(Y\).

\begin{align}
P(X = 0, Y = 0) &= P(\omega_8 = CCC) = 1/8; \notag\\
P(X = 1, Y = 0) &= P(\omega_5 = CCT) + P(\omega_6 = CTC) = 2/8; \notag\\
P(X = 1, Y = 1) &= P(\omega_7 = TCC) = 1/8; \notag\\
P(X = 2, Y = 0) &= P(\omega_4 = CTT) = 1/8; \notag\\
P(X = 2, Y = 1) &= P(\omega_3 = TCT) + P(\omega_2 = TTC) = 2/8; \notag\\
P(X = 3, Y = 1) &= P(\omega_1 = TTT) = 1/8; \notag
\end{align}

Le probabilità così trovate sono riportate nella tabella \ref{tab:ditr-cong-biv-1} la quale descrive la distribuzione di probabilità congiunta delle variabili aleatorie \(X\) = ``numero di realizzazioni con il risultato testa nei tre lanci'' e \(Y\) = ``numero di realizzazioni con il risultato testa nel primo lancio'' per l'esperimento casuale consistente nel lancio di tre monete equilibrate.

\begin{longtable}[]{@{}ccc@{}}
\caption{\label{tab:ditr-cong-biv-1} Distribuzione di probabilità congiunta per i risultati dell'esperimento consistente nel lancio di tre monete equilibrate.}\tabularnewline
\toprule
\(x / y\) & 0 & 1 \\
\midrule
\endfirsthead
\toprule
\(x / y\) & 0 & 1 \\
\midrule
\endhead
0 & 1/8 & 0 \\
1 & 2/8 & 1/8 \\
2 & 1/8 & 2/8 \\
3 & 0 & 1/8 \\
\bottomrule
\end{longtable}

In generale, possiamo dire che, dato uno spazio campionario discreto \(\Omega\), è possibile associare ad ogni evento elementare \(\omega_i\) dello spazio campionario una coppia di numeri reali \((x, y)\), essendo \(x = X(\omega)\) e \(y = Y(\omega)\), il che ci conduce alla seguente definizione.

\begin{definition}
Siano \(X\) e \(Y\) due variabili casuale. La funzione che associa ad ogni coppia \((x, y)\) un livello di probabilità prende il nome di funzione di probabilità congiunta:

\[
P(x, y) = P(X = x, Y = y).
\]
\end{definition}

\noindent
Il termine ``congiunta'' deriva dal fatto che questa probabilità è legata al verificarsi di una coppia di valori, il primo associato alla variabile casuale \(X\) ed il secondo alla variabile casuale \(Y\). Nel caso di due sole variabili casuali si parla di distribuzione bivariata, mentre nel caso di più variabili casuali si parla di distribuzione multivariata.

\hypertarget{proprietuxe0}{%
\subsection{Proprietà}\label{proprietuxe0}}

Una distribuzione di massa di probabilità congiunta bivariata deve soddisfare due proprietà:

\begin{enumerate}
\def\labelenumi{\arabic{enumi}.}
\item
  \(0 \leq P(x_i, y_j) \leq 1\);
\item
  la probabilità totale deve essere uguale a \(1.0\). Tale proprietà può essere espressa nel modo seguente
\end{enumerate}

\[
\sum_{i} \sum_{j} P(x_i, y_j) = 1.0.
\]

\hypertarget{eventi}{%
\subsection{Eventi}\label{eventi}}

Si noti che dalla probabilità congiunta possiamo calcolare la
probabilità di qualsiasi evento definito in base alle variabili
aleatorie \(X\) e \(Y\). Per capire come questo possa essere fatto,
consideriamo nuovamente l'esperimento casuale discusso in precedenza.

\begin{example}
Per la distribuzione di massa di probabilità congiunta riportata nella tabella~precedente si trovi la probabilità dell'evento \(X+Y \leq 1\).

Per trovare la probabilità richiesta dobbiamo semplicemente sommare le probabilità associate a tutte le coppie \((x,y)\) che soddisfano la condizione \(X+Y \leq 1\), ovvero

\begin{equation}
P_{XY}(X+Y \leq 1) = P_{XY}(0, 0) + P_{XY}(1, 0)= 3/8.\notag
\end{equation}
\end{example}

\hypertarget{regola-della-catena}{%
\subsection{Regola della catena}\label{regola-della-catena}}

Regola della catena permette il calcolo di qualsiasi membro della distribuzione congiunta di un insieme di variabili casuali utilizzando solo le probabilità condizionate.

\begin{definition}
Dati due eventi \(A\) e \(B\), la regola della catena (\emph{chain rule}) afferma che

\[
P(A \cap B) = P(A)P(B \mid A).
\]
\end{definition}

Nel caso di 4 eventi, per esempio, la regola della catena diventa

\[
P(A_1, A_2, A_3, A_4) = P(A_1) P(A_2 \mid A_1) P(A_3 \mid A_1, A_2) P(A_4 \mid A_1, A_2, A_3).
\]

\hypertarget{funzioni-di-probabilituxe0-marginali}{%
\subsection{Funzioni di probabilità marginali}\label{funzioni-di-probabilituxe0-marginali}}

Data la funzione di probabilità congiunta \(p(x, y)\) è possibile pervenire alla costruzione della funzione di probabilità della singola variabile casuale \(X\) o della singola variabile casuale \(Y\)

\begin{equation}
P(X) = \sum_y Pr(X, Y = y),
\end{equation}

\begin{equation}
P(Y) = \sum_x P(X = x, Y),
\end{equation}

\noindent
che prendono, rispettivamente, il nome di funzione di probabilità marginale di \(X\) e funzione di probabilità marginale di \(Y\). Si noti che \(P_X\) e \(P_Y\) sono normalizzate:

\[
\sum_x P_X(x) = 1.0, \quad \sum_y P_Y(y) = 1.0.
\]

\begin{example}
\protect\hypertarget{exm:flip3coins}{}\label{exm:flip3coins}

Per l'esperimento casuale consistente nel lancio di tre monete equilibrate, si calcolino le probabilità marginali di \(X\) e \(Y\).

Nell'ultima colonna a destra e nell'ultima riga in basso della tabella \ref{tab:ditr-cong-biv} sono riportate le distribuzioni di probabilità marginali di \(X\) e \(Y\). \(P_X\) si ottiene sommando su ciascuna riga fissata la colonna \(j\), \(P_X(X = j) = \sum_y p_{xy}(x = j, y)\). \(P_Y\) si trova sommando su ciascuna colonna fissata la riga \(i,\) \(P_Y (Y = i) = \sum_x p_{xy}(x, y = i)\).

\begin{longtable}[]{@{}cccc@{}}
\caption{\label{tab:ditr-cong-biv} Distribuzione di probabilità congiunta \(p(x,y)\) per i risultati dell'esperimento consistente nel lancio di tre monete equilibrate e probabilità marginali \(P(x)\) e \(P(y)\).}\tabularnewline
\toprule
\(x / y\) & 0 & 1 & \(P(x)\) \\
\midrule
\endfirsthead
\toprule
\(x / y\) & 0 & 1 & \(P(x)\) \\
\midrule
\endhead
0 & 1/8 & 0 & 1/8 \\
1 & 2/8 & 1/8 & 3/8 \\
2 & 1/8 & 2/8 & 3/8 \\
3 & 0 & 1/8 & 1/8 \\
\(P(y)\) & 4/8 & 4/8 & 1.0 \\
\bottomrule
\end{longtable}

\end{example}

\hypertarget{marginalizzazione-nel-caso-continuo}{%
\section{Marginalizzazione nel caso continuo}\label{marginalizzazione-nel-caso-continuo}}

Come indicato nel Paragrafo precedente, la marginalizzazione è un metodo che richiede di sommare tutti i possibili valori di una variabile casuale per determinare il contributo marginale di un'altra. Per una v.c. discreta abbiamo

\begin{equation}
Pr(X) = \sum_y Pr(X, Y = y),\notag
\end{equation}

\noindent
ovvero

\begin{equation}
P(X) = \sum_y P(X, Y = y) = \sum_y P(X \mid Y = y) P(Y = y).
\end{equation}

Nel caso continuo abbiamo

\begin{equation}
p(X) = \int_y p(X, Y = y) dy = \int_y p(X \mid Y = y) p(Y = y).
\end{equation}

\hypertarget{verosimiglianza-marginalizzata}{%
\section{Verosimiglianza marginalizzata}\label{verosimiglianza-marginalizzata}}

Nella statistica bayesiana, la marginalizzazione serve per calcolare la cosiddetta \emph{verosimiglianza marginalizzata}. Sia \(y = \{y_1, \dots, y_n\}\) un campione casuale e si assuma un modello statistico parametrizzato da \(\theta\). Il teorema di Bayes è

\[
p(\theta \mid y) = \frac{p(y \mid \theta) p(\theta)}{p(y)},
\]

\noindent
laddove

\[
p(y) = \int_\theta p(y \mid \theta) p(\theta) d\theta,
\]

\noindent
si integra su tutti i valori del parametro (o dei parametri).

\hypertarget{indipendenza-stocastica}{%
\section{Indipendenza stocastica}\label{indipendenza-stocastica}}

Ora abbiamo tutti gli strumenti per dare una precisa definizione
statistica al concetto di indipendenza. La definizione proposta
sarà necessariamente coerente con la definizione di indipendenza che abbiamo usato fino ad ora. Ma, espressa in questi nuovi termini, potrà essere utilizzata in indagini probabilistiche e statistiche più complesse. Ricordiamo che gli eventi \(A\) e \(B\) si dicono indipendenti se \(P (A \cap B)\, = P(A) P(B)\). Diciamo quindi che \(X\) e \(Y\) sono indipendenti se qualsiasi evento definito da \(X\) è indipendente da qualsiasi evento definito da \(Y\). La definizione formale che garantisce che ciò accada è la seguente.

\begin{definition}
Le variabili aleatorie \(X\) e \(Y\) sono indipendenti se la loro distribuzione congiunta è il prodotto delle rispettive distribuzioni marginali:

\begin{equation}
P(X, Y)\, = P_X(x)P_Y(y).
\end{equation}
\end{definition}

Nel caso discreto, dunque, l'indipendenza implica che la probabilità
riportata in ciascuna cella della tabella di probabilità congiunta deve
essere uguale al prodotto delle probabilità marginali di riga e di
colonna:
\[
P(x_i, y_i)\, = P_X(x_i) P_Y(y_i).
\]

\begin{example}
Per la situazione rappresentata nella tabella \ref{exm:flip3coins} le variabili casuali \(X\) e \(Y\) sono indipendenti?

Nella tabella le variabili casuali \(X\) e \(Y\) non sono indipendenti: le probabilità congiunte non sono ricavabili dal prodotto delle marginali. Per esempio, nessuna delle probabilità marginali è uguale a \(0\) per cui nessuno dei valori dentro la tabella (probabilità congiunte) che risulta essere uguale a \(0\) può essere il prodotto delle probabilità marginali.
\end{example}

\hypertarget{considerazioni-conclusive}{%
\section*{Considerazioni conclusive}\label{considerazioni-conclusive}}
\addcontentsline{toc}{section}{Considerazioni conclusive}

La funzione di probabilità congiunta tiene simultaneamente conto del
comportamento di due variabili casuali \(X\) e \(Y\) e di come esse si
influenzano reciprocamente. In particolare, si osserva che se le due
variabili non si influenzano, cioè se sono statisticamente indipendenti,
allora la distribuzione di massa di probabilità congiunta si ottiene
come prodotto delle funzioni di probabilità marginali di \(X\) e \(Y\):
\(P_{X, Y}(x, y) = P_X(x) P_Y(y)\).


% Bibliography
%%%%%%%%%%%%%%%%%%%%%%%%%%%%%%%%%%%%%%%%%%%%%%%%%%%%%%%%%%

\backmatter
\SmallMargins

\printbibliography
\onecolumn


% Tables (of tables, of figures)
%%%%%%%%%%%%%%%%%%%%%%%%%%%%%%%%%%%%%%%%%%%%%%%%%%%%%%%%%%


\cleardoublepage
\LargeMargins
\listoffigures


% After-body (LaTeX code inclusion)
%%%%%%%%%%%%%%%%%%%%%%%%%%%%%%%%%%%%%%%%%%%%%%%%%%%%%%%%%%




% Back cover
%%%%%%%%%%%%%%%%%%%%%%%%%%%%%%%%%%%%%%%%%%%%%%%%%%%%%%%%%%%

% Even page, small margins, no running head, no page number.
\evenpage
\SmallMargins
\thispagestyle{empty}

\begin{normalsize}

\begin{description}

\selectlanguage{italian}
\item[Abstract]
This document contains the material of the lessons of Psicometria B000286 (2021/2022) aimed at students of the first year of the Degree Course in Psychological Sciences and Techniques of the University of Florence, Italy.
\item[Keywords]
Data science, Bayesian statistics.
~\\

\end{description}

\end{normalsize}


\end{document}
